\section{Why the Profit Motive Is Useful and Insufficient}

\subsection{Profit as Signal, Not Proof}

Profit is sometimes presented as the opposite of a noble mission. Other times it is presented as proof that a business has served society. Neither view is reliable.

Imagine an educational organization that genuinely improves the capabilities of its students but loses money on every cohort. It may survive for a while through a benefactor, a grant, or accumulated savings. Unless it can secure a continuing subsidy, improve efficiency, or charge enough to cover its costs, it eventually has to reduce quality, shrink, or close. In this case profit is not the enemy of quality. Economic viability is 1 of the conditions for providing quality repeatedly.

Now imagine another educational company that sells career dreams to anxious students, makes implausible promises, spends heavily on sales, and delivers poor instruction. It may report attractive short-term margins. The profit does not establish that the company created corresponding human value. The business may simply have become effective at selling something difficult to evaluate before purchase.

Profit should therefore be treated as a feedback signal with real limitations. Sustained profits, earned under reasonably competitive and transparent conditions, often indicate that customers value what the company produces relative to its costs. But profits can also reflect temporary shortages, regulatory privileges, externalized costs, weak information, or the slow erosion of future quality.

\textbf{The useful founder's question is not whether profit is good or bad. It is whether the enterprise has discovered a repeatable and legitimate way to create enough value that customers willingly finance its continuation.}

This is the basis of sustainable value capture.

\subsection{When the Reward System Misleads}

It also explains why incentives matter. A tuition business rewarded only for enrollment will tend to devote energy to acquisition. One rewarded only for pass rates may weaken its tests. A consultancy paid exclusively for billable hours may have little incentive to eliminate the client's problem quickly. A salesperson paid only on signed contracts may promise things the delivery team cannot provide. A software founder rewarded only for user growth may postpone the awkward question of whether users will pay enough to support the product.

\textbf{Every business creates a system of rewards, and people respond to the things that system measures and pays for. An institution that hopes to endure must align its commercial survival with the customer outcomes it claims to value.} That alignment is difficult, incomplete, and continually revisable; it cannot be achieved by writing a mission statement alone.
